\chapter{Introduction}
\label{cha:intro}

% Multinomial data naturally arise in many applications where observations are distributed across a finite set of mutually exclusive categories. For example, in market research, the number of consumers choosing among several (competing) brands [CITATION]; in public health studies, the number of individuals falling into different health status categories [CITATION]; and in text analysis, the word count from a fixed vocabulary from sampled books [CITATION]. These multinomial counts are, however, often associated with covariates or grouping structures whose relationships with the observed counts are of substantive interest.

% Multinomial, common questions
Multinomial data arise naturally in many applications. Observations are distributed across a finite set of mutually exclusive categories. For example, \citet{zhang2026conjugatingvariationalinferencelarge} used mixed multinomial logit model to analyse the number of consumers choosing among several competing brands in market research. In public health, \citet{Beltran2024} models the number of individuals falling into different health status categories. In text analysis, word counts from a fixed vocabulary are used to characterise sampled documents \citep{Taddy_2015}. These multinomial counts rarely arise in isolation. They are typically associated with covariates describing the observations and group structure. The relationship between this structure and the observed counts is of substantive interest.

% Example, previous work: multinomial(pme, dmr) -> reduced-rank structutre 
A large body of previous work models the category probabilities as a function of covariates and group membership. The multinomial logistic regression model is the classical starting point. It extends naturally to a mixed-effects model that adds a group-level random effect to capture within-group correlation, which is an approach that still under active development at large scale for multinomial choice data with many groups \citep{zhang2026conjugatingvariationalinferencelarge}. In recent years, there has been a growing movement towards the application of high-dimensional multinomial datasets with large number of categories and groups. Previous work on large number of groups has been structured both theoretically and practically [Citation].

% Limitations in high-dim cases -> our work (fa-dmr)
However, the large-Q regimes are not discussed in the common work, which have arisen in many applications. As mentioned before, vocabulary sizes in text analysis are typically in the thousands \citep{Taddy_2015}. Taxon or gene counts in bioinformatics and genomics can run into the thousands. Spatial count data are often recorded over large numbers of small areal units \citep{Banerjee2008}. Due to the coupled nature of the multinomial distribution, computing the probability of a particular category requires considering the probabilities of all categories. This coupling not only increases the computationally cost as the number of categories grows, but also prevents the likelihood from being naturally decomposed across categories, posing challenges for scalable computation and inference. In the meanwhile, this bottleneck happens when we introduce the group structure $\bm{\lambda}$ into multinomial distribution, since the covariance $\operatorname{Cov}\left(\bm{\lambda}\right)$ has $\mathcal{O}(Q^2)$ free parameters. This is far more than any realistic number of groups $J$ can identify. It also means the $Q \times Q$ matrix operations that any likelihood-based procedure needs, such as inversion or computing a determinant, cost $O(Q^3)$, regardless of identifiability. Previous work faces obstructs on algorithms as well. \citet{Goplerud_2024} showed that standard mean-field coordinate ascent variational inference (CAVI) systematically under-estimates posterior uncertainty in high-dimensional mixed models with the degree of under-estimation worsening as dimension grows. \citet{Henderson02102019} faced similar difficulty for the alternative approach of maximum likelihood estimation via the EM algorithm, whose convergence rate is linear. This is a serious practical bottleneck for high-dimensional scenarios. 

The multinomial likelihood's row-total constraint $\sum_q \pi_{ijq}=1$ couples its categories together. This makes the observed-data Hessian dense and expensive to invert as $Q$ grows. 

% factor analyzer
A key approach for solving this obstruct is reduced-rank method. Multivariate random effects can be represented parsimoniously using a reduced-rank approach, which makes simplifying assumptions that reduce the dimension of the problem to $K\ll Q$ and fit a multivariate random effect via a linear combination of $K$ latent variables \citep{Bartholomew2011}. This framework is often referred to as a factor analytical model \citep{Bartholomew2011}, which has been used frequently in ecology and the social sciences. In the case of exponential family responses, it is sometimes called a generalized latent variable model (GLVM) \citep{Anders2004}.(check)
is not simply a matter of scale. Left unstructured, the group-level covariance $\Sigma = \mathrm{Cov}(\bm{\lambda})$ has $O(Q^2)$ free parameters. This is far more than any realistic number of groups $J$ can identify. It also means the $Q \times Q$ matrix operations that any likelihood-based procedure needs, such as inversion or computing a determinant, cost $O(Q^3)$, regardless of identifiability. Two broad families of estimation strategy are used to fit models of this kind. The first treats the group-level random effects $\bm{\lambda}_j$ as missing data, to be integrated out of the likelihood. This integration is done exactly where feasible, or via a Laplace approximation to the resulting intractable integral, within an EM algorithm that alternates between this approximate integration step and updating the remaining parameters. The second treats the same integration problem variationally, replacing the true posterior of $\bm{\lambda}$ with a tractable approximating family, most commonly a mean-field Gaussian fit by coordinate ascent (CAVI).

For multinomial distribution, many approaches are established to simplify the statistical modelling. \citet{Baker1994} exploited the Poisson-Multinomail Equivalence (PME) to decouple the multinomial likelihood as the product of independent-Poisson likelihood conditional on the individual total. \citet{Taddy_2015} adopted this work and established Distributed Multinomial Regression (DMR), which reduced a coupled $Q$-category regression to $Q$ independent Poisson regressions and applied on high-dimensional dataset.  
 
\todo{Describe what models are used to analyse this like below}

A natural approach to make inferences between covariates, cluster groups and the multionmial counts is to use a modelling approach. For $i = 1, 2, \dots, I$ in $j = 1, 2, ..., J$ and $q = 1, 2, ..., Q$, suppose that 

\begin{itemize}
\item  $y_{ijq}$ is the count in category $q$ for the $i$-th observation in cluster group $j$,
\item $\boldsymbol{y}_{ij} = (y_{ij1}, y_{ij2}, \cdots, y_{ijQ})^\top \in \mathbb{Z}^{Q}_{\geq 0}$ is the $Q\times 1$ vector of counts for the $i$-th observation and $j$-th group, and 
\item $x_{ijk}$ is the $k$-th covariate for the $i$-th observation in the $j$-th group.
\end{itemize}

\noindent Then we model ....

\todo{Introduce the problem when $Q$ is large}

with a massive number of categories have been used extensively in text analysis \citep{Taddy_2015}, bioinformatics \citep{Smith2001, Buettner2015} (socail science, economic) and spatial statistics \citep{Cressie2008, Banerjee2008}. 


\todo{Then describe the parameter estimation methods}

 A natural way to resolve this tension between flexibility and tractability is to impose a reduced-rank, factor-analytic structure directly on group-level random effect $\bm{\lambda}$. This avoids estimating an unstructured $\Sigma$ and then correcting for the resulting instability. Write $\alpha_j = B f_j + \varepsilon_j$, for a $Q \times K$ loading matrix $B$ with $K \ll Q$, a $K$-dimensional latent factor score $f_j$, and idiosyncratic noise $\varepsilon_j \sim \mathcal{N}(0, \sigma^2 I_Q)$. This reduces the number of free covariance parameters from $O(Q^2)$ to $O(QK)$ \citep{Bartholomew2011}. This construction is the classical factor-analytic model \citep{AndersonRubin1956,Bartholomew2011}, long used in ecology and the social sciences. It is sometimes called a generalised latent variable model (GLVM) when the observed responses follow an exponential-family distribution \citep{Anders2004}\todo{verify this GLVM attribution before submission}. A related, parallel line of work reaches for the same reduced-rank idea from a log-normal-Poisson rather than a multinomial starting point \citep{Batardiere2024}. We discuss the relationship between the two traditions in Chapter~\ref{cha:model-specification}.
 
In this thesis, we start from data that follow the multinomial distribution directly, which is the natural distribution for count data rather than one instance of a general exponential-family response. We denote latent factors $\mathbf{f}$ as missing data and $\bm{\lambda}$ as group-level random effects and imposes the factor-analytic and reduced-rank structure discussed above on the decomposition $\bm{\lambda}=\mathbf{B}\mathbf{B}^{\top}+\mathbf{D}$, then use an Expectation-Maximization (EM) algorithm to handle the resulting incomplete-data problem with unobserved $\mathbf{f}$ . Using the Poisson-multinomial equivalence introduced above \citep{Baker1994, Taddy_2015, lee2017poissontrickextensionsfitting}, we decouple thecomputation of Hessian. This accelerates each EM iteration from cubic to near-linear cost in $Q$. We further establish conditions under which the resulting estimator is consistent and asymptotically normal. These conditions are expressed in terms of the joint growth of the number of groups $J$ and the number of categories $Q$, rather than in terms of $Q$ alone. Recent work has pursued a related reduced-rank idea for item-level count data via joint maximum likelihood \citep{Cui2026} instead of EM algorithm. They discussed EM algorithm and its statistical properties as an open problem since the marginal likelihood maximization suffers more difficulties when dealing with integration of latent factors. As well,  we extend their model structure particularly in allowing genuine idiosyncratic variance to be considered, rather than the reduced-rank structure only. 
 
The remainder of the thesis is organised as follows. Chapter~\ref{cha:model-specification} introduces the model in full with identifiability conditions. It reviews the high-dimensional multivariate count factor-analytic literature and its extensions, and introduces the PME and the classical EM algorithm formally. Chapter~\ref{cha:methodology} develops the dense EM algorithm and the PME\&Woodbury-accelerated version. It discusses convergence and sources of approximation error, and compares the computational cost of our approach with the methods surveyed above. Chapter~\ref{cha:simulation-result} reports simulation studies. Chapter~\ref{cha:real-result} applies the model to real grouped count data. Chapter~\ref{cha:conc} concludes, discussing limitations and directions for future work. Commonly used notation is summarized in Table~\ref{tab:notation}.

\iffalse
In this thesis, we start with data that follows the Multinomial distribution, which is always used in counting dataset, instead of starting with the exponential family. Focus on reduced-rank structure in factor-analytic model, we use Expectation-Maximization (EM) Algorithm to deal with incomplete data with unknown latent variables. Given the Poisson-Multinomial Equivalence (PME) \citep{10.2307/2348134, lee2017poissontrickextensionsfitting}, we decouple the structure of Multinomial Hessian matrix and accelerate the EM algorithms with latent variables. 

In the next Chapter, we provide our model setup with identifiability conditions, an overview of the (high-dimensional) multivariate count factor-analytic models and extensions, the introduction of PME, and the brief introduction of classical EM algorithm. In Chapter~\ref{cha:methodology}, we describe the dense EM algorithm and Woodbury acceleration method decoupling by PME auxiliary variable structure. We discuss its convergence and sources of error. Furthermore, we compare the performance and computational efficiency of our proposed model with those of competing models reported in the literature review. Simulation studies and real data analysis are given in Chapter~\ref{cha:simulation-result} and \ref{cha:real-result}. Chapter~\ref{cha:conc} concludes the thesis, giving the limitation and future work.  

This thesis is organised as follows. The notation commonly used is summarised in Table~\ref{tab:notation}. Chapter~\ref{cha:model-specification} describes the background,  Chapter~\ref{cha:methodology} describes the main methodology in this thesis, Chapter~\ref{cha:simulation-result} shows the simulation result, Chapter~\ref{} demonstrates a real data application and we end with Chapter.
\fi